\documentclass{article}
\usepackage{amsfonts,amsthm,amsmath,amssymb,mathtools}
\usepackage{mathrsfs}
\usepackage{enumitem}

\setcounter{section}{0}
\renewcommand{\thesection}{\S\arabic{section}}
\renewcommand{\thesubsection}{\arabic{section}}

\newtheorem{thm}{Theorem}
\newenvironment{theorem}[1]{
	\renewcommand{\thethm}{#1}
	\begin{thm}
		}{\end{thm}}

\newtheorem{lma}{Lemma}
\newenvironment{lemma}[1]{
	\renewcommand{\thelma}{#1}
	\begin{lma}
		}{\end{lma}}

\newcounter{section-number}
\setcounter{section-number}{1}

\theoremstyle{definition}
\newtheorem{ex}{}[section-number]
\newenvironment{exercise}[1]{
  \renewcommand{\theex}{\arabic{section-number}.\arabic{ex}}
	\begin{ex}
		}{\end{ex}}

\title{Exercises I.1}
\author{Connor Mooneyhan}
\date{June 7, 2024}

\begin{document}

\maketitle

\begin{ex}
	Locate a discussion of Russell's paradox, and understand it.
\end{ex}
\begin{proof}
	Done!
\end{proof}

\begin{ex}
	Prove that if $\sim$ is an equivalence relation on a set $S$, then the corresponding family $\mathscr{P}_\sim$ defined in \S1.5 is indeed a partition of $S$: that is, its elements are nonempty, disjoint, and their union is $S$.
\end{ex}
\begin{proof}
	If $S = \emptyset$, then $\mathscr{P}_\sim = \emptyset$ due to there being no elements with respect to which one can form equivalence classes, and $\emptyset$ does (vacuously) satisfy the conditions for a partition.

	If $S$ is nonempty, then given an equivalence class $[a]_\sim \in \mathscr{P}_\sim$, $a \in [a]_\sim$ since $a \sim a$.

	For disjointness, let $b \in S$ such that $[b]_\sim \neq [a]_\sim$.
	We proceed to show that $b$ is not related to $a$.
	Assume that $b \sim a$, and thus $a \sim b$.
	Then given $x \in [a]_\sim$, $x \sim a$, so by transitivity \begin{equation*}
		x \sim a \sim b \implies x \sim b,
	\end{equation*}
	so $x \in [b]_\sim$.
	A similar argument follows for any $y \in [b]_\sim$, so we conclude that $[a]_\sim = [b]_\sim$ despite the assumption to the contrary.
	Thus $b$ is not related to $a$.

	Given $x \in [a]_\sim$, by definition $x \sim a$, so $a \sim x$ by symmetry.
	If $x \in [b]_\sim$, then $x \sim b$, so the transitivity of $\sim$ gives us \begin{equation*}
		a \sim x \sim b \implies a \sim b,
	\end{equation*}
	contradicting our assumption.
	A similar argument can be made by swapping $a$ and $b$, so we conclude that $[a]_\sim$ and $[b]_\sim$ are disjoint.

	Finally, we show that $\bigcup \mathscr{P}_\sim = S$.
	Given $x \in \bigcup \mathscr{P}_\sim$, $x \in [a]_\sim$ for some $a \in S$, so by definition $x \in S$ since $[a]_\sim \subset S$.
	Assume, then, that $y \in S$.
	Then $y \in [y]_\sim \in \mathscr{P}_\sim$, we conclude that $y \in \bigcup \mathscr{P}_\sim$.
\end{proof}

\begin{ex}
	Given a partition $\mathscr{P}$ on a set $S$, show how to define a relation $\sim$ on $S$ such that $\mathscr{P}$ is the corresponding partition.
\end{ex}
\begin{proof}
	Given $x \in S$, let $P_x$ be the unique $P_x \in \mathscr{P}$ such that $x \in P_x$.
	Define $\sim$ such that for any $a, b \in S$, $a \sim b \iff P_a = P_b$.
	We first show that this is an equivalence relation.

	Given $a \in S$, clearly $P_a = P_a$, so $a \sim a$.
	Also, if $b \in S$ and $a \sim b$, then \begin{equation*}
		P_a = P_b \implies P_b = P_a \implies b \sim a.
	\end{equation*}
	Finally, if $c \in S$ and $b \sim c$, then \begin{equation*}
		P_a = P_b = P_c \implies P_a = P_c.
	\end{equation*}

	Now we show that $S/{\sim} = \mathscr{P}$.
	Let $[a]_\sim \in S/{\sim}$.
	We want to show that $[a]_\sim = P_a \in \mathscr{P}$, where $P_a$ is defined as above.
	Given $x \in [a]_\sim$, $x \sim a$, so by definition of $\sim$, $P_x = P_a$.
	Thus since $x \in P_x$, $x \in P_a$.
	On the other hand, if $y \in P_a$, then since $y$ is in only one member of $\mathscr{P}$ and $y \in P_y$, $P_y = P_a$ and thus $y \sim a \implies y \in [a]_\sim$.

	Let $A \in \mathscr{P}$.
	Then $A = P_b$ for some $b$ since $A$ is nonempty.
	We want to show that $P_b = [b]_\sim \in S/{\sim}$.
	Given $x \in P_b$, since $x \in P_x$ and $x$ is in only one member of $\mathscr{P}$, $P_x = P_b$.
	Thus $x \sim b$, so $x \in [b]_\sim$.
	If $y \in [b]_\sim$, then $P_y = P_b$, so since $y \in P_y$, $y \in P_b$.
\end{proof}

\begin{ex}
	How many different equivalence relations may be defined on the set $\left\lbrace 1, 2, 3 \right\rbrace$?
\end{ex}
\begin{proof}
	Five.
	This can be found easily by simply determining the number of partitions.
	Here's the list: \begin{align*}
		 & \left\lbrace \left\lbrace 1 \right\rbrace, \left\lbrace 2 \right\rbrace, \left\lbrace 3 \right\rbrace \right\rbrace \\
		 & \left\lbrace \left\lbrace 1 \right\rbrace, \left\lbrace 2, 3 \right\rbrace \right\rbrace                            \\
		 & \left\lbrace \left\lbrace 1, 3 \right\rbrace, \left\lbrace 2 \right\rbrace \right\rbrace                            \\
		 & \left\lbrace \left\lbrace 1, 2 \right\rbrace, \left\lbrace 3 \right\rbrace \right\rbrace                            \\
		 & \left\lbrace \left\lbrace 1, 2, 3 \right\rbrace \right\rbrace                                                       \\
	\end{align*}
\end{proof}

\begin{ex}
	Give an example of a relation that is reflexive and symmetric but not transitive.
	What happens if you attempt to use this relation to define a partition on the set?
	(Hint: Thinking about the second question will help you answer the first one.)
\end{ex}
\begin{proof}
	Let $S = \left\lbrace 1, 2, 3 \right\rbrace$, and define $\sim$ by \begin{align*}
		 & 1 \sim 1  \\
		 & 1 \sim 2  \\
		 & 2 \sim 1  \\
		 & 2 \sim 2  \\
		 & 2 \sim 3  \\
		 & 3 \sim 2  \\
		 & 3 \sim 3.
	\end{align*}
	It is easily checked that $\sim$ is reflexive and symmetric but not transitive.
	If you attempt to define a partition on $S$ using $\sim$ in the usual way, then you get \begin{equation*}
		\left\lbrace [1]_\sim, [2]_\sim, [3]_\sim \right\rbrace = \left\lbrace \left\lbrace 1, 2 \right\rbrace, \left\lbrace 1, 2, 3 \right\rbrace, \left\lbrace 2, 3 \right\rbrace \right\rbrace,
	\end{equation*}
	which is not a partition since its members are not disjoint.
\end{proof}

\begin{ex}
	Define a relation $\sim$ on the set $\mathbb{R}$ of real numbers by setting $a \sim b \iff b - a \in \mathbb{Z}$.
	Prove that this is an equivalence relation, and find a 'compelling' description for $\mathbb{R}/{\sim}$.
	Do the same for the relation $\approx$ on the plane $\mathbb{R} \times \mathbb{R}$ defined by declaring $(a_1, a_2) \approx (b_1, b_2) \iff b_1 - a_1 \in \mathbb{Z}$ and $b_2 - a_2 \in \mathbb{Z}$.
\end{ex}
\begin{proof}
	Given $a \in \mathbb{R}$, $a - a = 0 \in \mathbb{Z}$, so $a \sim a$.
	If $a \sim b$ for some $b \in \mathbb{R}$, then $b - a \in \mathbb{Z}$, so $a - b = -(b - a) \in \mathbb{Z}$ and thus $b \sim a$.
	Finally, if $c \in \mathbb{R}$ such that $b \sim c$, then $c - a = c - b + b - a \in \mathbb{Z}$ since both $c - b$ and $b - a$ are in $\mathbb{Z}$.

  As for a compelling description, you could characterize each member of $\mathbb{R}/{\sim}$ as simply $\mathbb{Z}$ translated by some $x \in [0, 1)$.

	Now we shift our focus to $\approx$, where we aim to show not only that $\approx$ is an equivalence relation on its own, but that it \textit{inherits} the properties of reflexivity, symmetry, and transitivity from $\sim$.
	To do this, we notice that we can recharacterize the definition of $\approx$ as \begin{equation*}
		(a, b) \approx (c, d) \iff a \sim c\ \mathrm{and}\ b \sim d.
	\end{equation*}
  This is clear from the definitions of $\approx$ and $\sim$.

  First, we notice that reflexivity follows from the fact that given $(a, b) \in \mathbb{R}^2$, $a \sim a$ and $b \sim b$, so $(a, b) \approx (a, b)$.
  Then assume that $(a, b) \approx (c, d)$ for some $(c, d) \in \mathbb{R}^2$.
  Then $a \sim c$ and $b \sim d$ by our re-definition, so $c \sim a$ and $d \sim b$, implying that $(c, d) \approx (a, b)$.
  Finally, given $(e, f) \in \mathbb{R}^2$ such that $(c, d) \approx (e, f)$, we have \begin{equation*}
    a \sim c \sim e \implies a \sim e
  \end{equation*}
  and \begin{equation*}
    b \sim d \sim f \implies b \sim f,
  \end{equation*}
  so $(a, b) \approx (e, f)$.

  We carry out the proof in this way so that it is apparent that nothing in the proof depends on the particulars of $\sim$.
  This hints at the appropriateness of characterizing $\approx$ as an instance of a "product relation", namely \begin{equation*}
    \approx = \sim \times \sim.
  \end{equation*}
  Our proof is sufficient to show that any such product of an equivalence relation with itself is itself an equivalence relation.

  Enough generalization.
  A compelling description for $S/{\approx}$ is similar to that of $S/{\sim}$, where for each member of $S/{\approx}$, we imagine the lattice $\mathbb{Z}^2$ translated by $(x, y) \in [0, 1) \times [0, 1)$.
\end{proof}
\end{document}
